\documentclass[11pt]{article}
\usepackage[a4paper,margin=1in]{geometry}
\usepackage{amsmath,amssymb,hyperref}
\title{The ratio in Conjecture 00000005135 has the wrong upper endpoint}
\author{gaochengzhecpu}
\date{October 3, 2026}
\hypersetup{hidelinks}
\begin{document}
\maketitle

The source explicitly takes the ratio of the \emph{structured} condition
number to the \emph{unstructured} one, and asserts an attainable upper
endpoint $\sqrt n$. Under the standard comparison, which retains the same
norm and sensitivity functional while restricting the perturbations, every
defined finite positive-denominator ratio is at most one. The proposed
endpoint already fails at $n=4$.

Write $\mathcal H_n$ for the vector space of Hankel matrices:
\[
 E_{ij}=E_{k\ell}\quad\text{whenever }i+j=k+\ell.
\]
Let $a(E)$ be the amplification associated with a perturbation direction.
For a differentiable matrix function it is the norm of its derivative
applied to $E$. The relevant absolute condition numbers are
\[
 \kappa_{\rm u}=\sup_{\|E\|\leq1}a(E),\qquad
 \kappa_{\rm s}=\sup_{\substack{\|E\|\leq1\\E\in\mathcal H_n}}a(E).
\]
Any common positive normalization for relative condition numbers cancels
in their ratio. The structured domain is a subset of the unstructured
domain. The defining least-upper-bound property therefore gives
\[
 \kappa_{\rm s}\leq\kappa_{\rm u}.
\]
This standard structured/unstructured comparison is also explicitly stated
by Arslan, Noferini, and Tisseur\cite{ant}, who interpret a small ratio
in this direction as a benefit of preserving structure.

Whenever $0<\kappa_{\rm u}<\infty$ and the ratio is defined,
\[
 r=\frac{\kappa_{\rm s}}{\kappa_{\rm u}}\leq1.
\]
Consequently the supremum over \emph{any} family of these ratios, when
it exists as a finite least upper bound, is also at most one. In dimension
four the asserted upper endpoint is $\sqrt4=2>1$, so it cannot be attained,
either by an individual pair of condition numbers or by a supremum over
symbol pairs. The upper-end attainment clause of the conjecture is false.

The positivity condition does not discard a possible endpoint witness:
a finite ordinary ratio equal to two requires a nonzero denominator.
Zero-over-zero and infinite-over-infinite do not define such a real ratio.
This argument applies to the standard common-norm comparison. Comparing
a parameter-vector norm on one side with a differently scaled matrix norm
on the other would be a different definition; no such change is specified
in the source.

\newpage
\paragraph{Formal verification.}
The Lean 4.19.0 file \texttt{lean/Main.lean} uses only bundled Std.
It defines Hankel matrices by all anti-diagonal equalities and defines
both condition numbers by the full least-upper-bound predicate on their
actual perturbation domains. The amplification and norm functions are
arbitrary: proving the statement for all such functions covers every
choice arising from a matrix derivative.

The elementary scalar laws are packaged as the universally quantified
\texttt{OrderedScale} interface. They are transitivity, positive
multiplication cancellation, the multiplicative unit law, and ordinary
strict-order facts. They are satisfied by real scalars; an explicit
integer model is also included to verify consistency of this interface.
The result is not restricted to integer-valued condition numbers.
The equation $r\kappa_{\rm u}=\kappa_{\rm s}$ represents a ratio without
division. The final theorem excludes a supremum of two for \emph{every}
family of admissible Hankel ratios in dimension four, using the actual
leastness property of the supremum. It is named
\begin{quote}\texttt{conjecture\_5135\_endpoint\_false}.\end{quote}
Its printed axiom list is empty. No \texttt{sorry}, \texttt{admit},
\texttt{native\_decide}, or extra axiom is used.

\begin{thebibliography}{1}
\bibitem{ant}
B. Arslan, V. Noferini, and F. Tisseur,
\emph{The Structured Condition Number of a Differentiable Map Between
Matrix Manifolds, with Applications}, MIMS EPrint 2017.36,
especially the comparison on page 2 and Section 2.
\url{https://eprints.maths.manchester.ac.uk/2579/}.
\end{thebibliography}
\end{document}
